\chapter{まとめと今後の検討事項}
\label{ch:conclusion}

\section{本稿のまとめ}

本稿では，チーム 23 のテーマ「Arduino オーケストラ」について，筆者が担当する指揮者マイコン部を中心に以下を整理した。

\begin{itemize}
  \item 合奏への参加障壁低減と STEAM 教育への接続という \textbf{社会的背景}
  \item 無線マイコン・内蔵 IMU・低遅延 UDP 通信・Processing 音響合成という \textbf{技術的背景}
  \item Wii Music・バーチャル指揮者研究・ESP32 分散演奏・JackTrip など \textbf{既存システム}の位置付け
  \item IMU ジェスチャ $\rightarrow$ 演奏パラメータの写像と，UDP 分散同期を組み合わせる \textbf{本提案のアイディア}
\end{itemize}

既存システムと比較して，「物理的に分散した楽器ノードを人間の指揮動作でリアルタイム制御する」点が本提案の差別化点である。

\section{チームで協議したい事項}

次回ミーティング（2026-04-22）では，以下を中心に協議したい。

\begin{enumerate}
  \item \textbf{課題曲の選定}: 4 パート構成に分解しやすい輪唱曲（森のクマさん／かえるのうた／静かな湖畔 など）の比較。
  \item \textbf{パケット仕様の詳細化}: BEAT / CTRL パケットのフィールド構成・送信頻度・ロス対策。
  \item \textbf{指揮ジェスチャの写像の妥当性}: 振幅$\rightarrow$強弱，高周波$\rightarrow$ビブラートの対応が直感的か，実機で検証できるか。
  \item \textbf{ストレッチ目標の優先度}: 強弱・ビブラート・視覚演出のうち，限られた開発期間でどこまで狙うか。
  \item \textbf{評価方法}: 同期精度（タイミング誤差）や楽器音らしさの定量・定性評価の方法。
\end{enumerate}

\section{今後の個人タスク}

次回までに筆者が進める作業は以下である。

\begin{itemize}
  \item UNO R4 WiFi 内蔵 IMU からの加速度・角速度取得プロトタイプの作成
  \item 加速度ノルムのピーク検出による拍抽出アルゴリズムの試作
  \item UDP ブロードキャスト送信のサンプル実装と遅延計測
\end{itemize}

これらを通じて，次章以降の設計書（ハッカソン2 向け）で詳細化するための実データ・数値根拠を収集する。
